\chapter{Discussion}

This chapter interprets the results of the previous chapter without introducing any new figures. The objective is to read the evidence honestly: what the pipeline demonstrates, what it does not, and what the reinforcement learning formulation actually provides compared to a strong supervised baseline.

\section{In-Distribution versus Hard Generalization}
The most informative contrast in the work is not the in-distribution figure, but its distance from the day-split evaluation. On the fixed random split, the MAIN model obtains an accuracy of 0.99381; under the day split (Check C), it drops to 0.84135 and the attack recall falls from 0.99536 to 0.52954. This gap should not be read as a failure, but as the honest measure of the problem: the random split shares the distribution between training and testing and, as quantified by the duplicate analysis (24.86\% of the test rows duplicate a training row, and up to 40.12\% among the test attacks), a considerable part of its performance is memorization. The day split breaks that comfort and reveals the true difficulty of generalizing to different temporal patterns and attack families. Reporting both rungs, rather than only the most favorable one, is the core methodological contribution of this thesis.

\section{QRDQN versus the Supervised Baseline}
The same-protocol comparison allows for a nuanced and deliberately non-triumphalist conclusion. \emph{In-distribution}, the balanced Random Forest baseline matches---and even slightly exceeds---QRDQN (attack F1 0.99676 versus 0.98445). Therefore, this thesis does not claim that reinforcement learning is superior on the random rung; on tabular in-distribution flow data, a well-configured Random Forest is a very strong competitor, consistent with the literature \parencite{LiuLang2019IDSSurvey,Apruzzese2022CrossEvaluationNIDS}.

The difference appears under distribution shift. On the same day split (Monday/Tuesday/Wednesday\,$\rightarrow$\,Thursday/Friday), both models degrade, but in very different ways: the Random Forest retains high attack precision (0.96473) at the expense of an attack recall that collapses to 0.08135---meaning it lets more than 91\% of attacks pass---while QRDQN maintains an attack recall of 0.52954. In terms of operational security, where a false negative is the expensive error, the cost-sensitive agent retains a much larger fraction of detection capability under adverse conditions. This divergence is attributable to the reward asymmetry ($\mathrm{fn}=-5.0$ vs $\mathrm{fp}=-2.0$), which pushes the policy toward blocking uncertain attacks rather than towards aggregate precision. The honest reading is not "RL wins", but rather that \emph{the cost-sensitive formulation degrades more favorably on the critical security metric under domain shift}, even though both approaches fall well below their in-distribution performance and neither is ready for deployment.

\section{Phase 2 Behavior and Domain Shift}
The Phase 2 inference offers a third perspective. On operator-generated laboratory traffic, the MAIN model achieves an accuracy of 0.991862 with benign distribution diagnostics (max absolute z 9.547, mean 0.3417, no clipping activated). This indicates that the laboratory capture is close to the training distribution and that the robust preprocessing pipeline---percentile clipping and z-score clipping on the persisted scaler---transfers correctly. However, this result must be interpreted with caution: unlike the day split, which subjects the model to an \emph{adverse} distribution shift within CICIDS2017, the laboratory capture is proprietary, non-adversarial traffic from a single closed environment. A strong Phase 2 result is, therefore, a controlled check of limited external validity and not proof of robustness against real-world domain shift; the two generalization signals---Check C and Phase 2---are complementary and must not be confused.
